El análisis exploratorio permitirá describir la cobertura anual, el número de partidos y jugadoras disponibles, la distribución de minutos y posesiones, así como la frecuencia de cada quinteto y de sus enfrentamientos observados. Se estudiarán los valores extremos y la variabilidad de los ratings ofensivo, defensivo y neto para identificar estimaciones inestables por tamaño de muestra.

También se analizarán distribuciones y correlaciones de puntos, asistencias, rebotes, pérdidas, robos, bloqueos y métricas de tiro, con el propósito de seleccionar atributos adecuados para caracterizar perfiles. Las visualizaciones y resúmenes se presentarán sólo después de ejecutar la exploración; no se reportan cifras observadas en esta propuesta.
Como el conjunto de datos esta presente en formato parquet se hará uso de la herramienta previamente mencionada de visual studio code llamada Parquet Visualizer.
\begin{lstlisting}[
    style=sqlstyle,
    caption={Lectura e integración de archivos play-by-play por temporada},
    label={lst:integracion_pbp}
]
SELECT
    *,
    regexp_extract(
        filename,
        '(\d{4})\.parquet$',
        1
    ) AS anio_doc
FROM read_parquet(
    'd:/Progra/Escom/Trabajo terminal/codigos_trabajo_terminal/datos_a_usar/wnba_stats/pbp/parquet/wnba_play_by_play_*.parquet',
    filename = true
);
\end{lstlisting}
Dentro de esta herramienta podremos juntar todos los tipos de
Se realizo un análisis exploratorio inicial de las variables presentes preliminar del conjunto de datos que se cuenta.
Posteriormente la herramienta permite exportar a csv que es el formato que ocuparemos para el eda.
\begin{lstlisting}[
    style=pythonstyle,
    caption={Lectura de datos play-by-play de la WNBA},
    label={lst:lectura_pbp}
]
import pandas as pd
import numpy as np
import matplotlib.pyplot as plt

file = 'play_by_play_2026.csv'
df = pd.read_csv(file)
df.dtypes
\end{lstlisting}
El codigo \ref{lst:lectura_pbp} nos presenta una visualizacion del tipo de datos presente en el archivo play by play con la siguiente salida.
\begin{table}[H]
\centering
\caption{Estructura de datos y tipos de variables}
\label{tab:variables_wnba}
\begin{tabular}{ll}
\toprule
\textbf{Columna} & \textbf{Tipo de Dato} \\
\midrule
order\_index & int64 \\
action\_number & int64 \\
clock & str \\
period & int64 \\
team\_id & int64 \\
team\_tricode & str \\
person\_id & int64 \\
player\_name & str \\
player\_name\_i & str \\
x\_legacy & int64 \\
y\_legacy & int64 \\
shot\_distance & int64 \\
shot\_result & str \\
is\_field\_goal & int64 \\
score\_home & float64 \\
score\_away & float64 \\
points\_total & int64 \\
location & str \\
description & str \\
action\_type & str \\
sub\_type & str \\
video\_available & int64 \\
shot\_value & int64 \\
action\_id & int64 \\
game\_id & int64 \\
seconds\_remaining & float64 \\
event\_type & str \\
is\_made\_shot & bool \\
is\_missed\_shot & bool \\
is\_free\_throw & bool \\
is\_rebound & bool \\
is\_turnover & bool \\
is\_foul & bool \\
is\_substitution & bool \\
is\_jump\_ball & bool \\
is\_timeout & bool \\
is\_period & bool \\
season & int64 \\
filename & str \\
anio\_doc & int64 \\
\bottomrule
\end{tabular}
\end{table}
De manera analoga se pueden obtener estas metricas atravez de la herramienta mencionada llamada data wrangler con ella se lograron obtener los siguientes resultados donde se detalla el procedimiento a realizaar a cada conjunto de datos.